\chapter*{Введение}
\addcontentsline{toc}{chapter}{Введение}

Задачи траекторного управления в трехмерном пространстве возникают в беспилотной авиации, мобильной робототехнике и других системах, которые должны двигаться по заданной пространственной кривой с ограничениями по точности, скорости и устойчивости. На практике этого обычно недостаточно свести к движению по заранее известной траектории. Важно еще и подобрать такой режим движения, при котором объект проходит маршрут достаточно быстро, но при этом не выходит за допустимые ошибки и не теряет устойчивость.

Для нелинейных динамических систем эта задача особенно чувствительна к качеству модели и способу построения регулятора. Линейные приближения удобны на этапе первичного анализа, но при заметных изменениях режима движения, кривизны траектории или ограничений объекта их возможностей часто уже не хватает. Поэтому на первый план выходит сочетание аналитических методов управления, которые дают понятную структуру закона управления, и интеллектуальных методов, которые позволяют адаптировать отдельные параметры под текущую ситуацию.

В качестве математической основы в работе используется подход к согласованному траекторному управлению по выходу, изложенный в диссертации С.\,А.~Ким~\cite{kim_algorithms}. Этот подход был выбран как базовый, поскольку он позволяет формализовать задачу слежения за пространственной кривой, ввести ошибки в сопутствующей системе координат и построить регулятор, учитывающий структуру нелинейной модели объекта.

Проблема, на решение которой направлена работа, состоит в следующем. Базовый регулятор обеспечивает устойчивое движение вдоль траектории, однако при фиксированном значении параметрической скорости $V^{\ast}$ система работает слишком консервативно: на простых участках траектории объект движется медленнее, чем мог бы, а на более сложных участках запас по устойчивости, наоборот, быстро уменьшается. Из-за этого возникает необходимость в адаптивном выборе скорости движения, согласованном с текущим состоянием объекта и геометрией траектории.

Цель работы состоит в исследовании методов интеллектуального управления траекторным движением для класса нелинейных систем на примере квадрокоптера и в разработке программной реализации, в которой аналитический регулятор траекторного слежения дополняется интеллектуальным модулем выбора параметрической скорости.

Для достижения поставленной цели были сформулированы следующие задачи:
\begin{enumerate}
    \item рассмотреть постановку задачи траекторного управления нелинейным объектом в трехмерном пространстве;
    \item описать математическую модель квадрокоптера и ограничения, существенные для задачи слежения;
    \item реализовать базовый алгоритм траекторного слежения по выходу;
    \item подготовить процедуру формирования датасета для подбора допустимой скорости движения;
    \item разработать и обучить нейросетевой модуль, выбирающий значение $V^{\ast}$ по текущим признакам состояния;
    \item провести численное моделирование и сравнить базовый и адаптивный режимы управления.
\end{enumerate}

Практическая значимость работы заключается в том, что полученный программный комплекс можно использовать как исследовательский стенд для проверки алгоритмов траекторного управления и для сравнения различных способов адаптивного выбора скорости движения. Кроме того, сама архитектура решения показывает, как можно сочетать аналитический регулятор, отвечающий за устойчивость, и интеллектуальную надстройку, отвечающую за более эффективный режим движения.
